%% 第五章

\chapter{总结与展望}
\chaptermark{总结与展望}

\section{全文总结}

本文围绕大语言模型偏好对齐中的奖励建模稳定性问题展开研究，提出了一种面向偏好对齐的近端逆奖励优化方法PIRO。针对传统RLHF范式中固定奖励模型容易受到策略分布偏移影响、进而引发奖励过优化和训练不稳定的问题，本文从逆强化学习角度出发，将偏好对齐过程建模为奖励函数学习与策略优化交替进行的动态过程，使奖励函数能够随策略生成分布的变化而更新。

在理论层面，本文将语言模型偏好对齐形式化为奖励学习与策略优化耦合的IRL双层优化问题。内层在给定奖励函数的条件下进行带KL正则的策略优化，外层则通过高质量示范响应学习奖励函数。基于该形式化，本文分析了奖励函数、KL正则强度与策略偏移之间的关系，并构造了固定旧策略的局部替代目标。理论分析表明，该替代目标在当前展开点处与原始目标具有相同的函数值和一阶梯度信息，能够在局部小步更新范围内提供一致的优化方向。进一步地，本文给出了带误差项的改进下界：
\begin{equation}
    L(\theta_{\mathrm{new}}) \ge L_{\mathrm{old}}(\theta_{\mathrm{new}})-C\epsilon .
\end{equation}
该结果说明，当相邻轮次奖励函数变化幅度受到约束时，替代目标相对于原始目标的偏差可以被显式控制，从而为受约束动态奖励更新提供了理论依据。

在算法层面，本文将理论中的奖励变化控制转化为可计算的RMSD近端约束，并设计自适应约束系数调节机制。RMSD用于刻画相邻轮次奖励函数在训练批次上的变化幅度，自适应约束系数则根据奖励变化是否超过目标范围动态调整惩罚强度。当奖励函数更新过快时，算法增强近端约束以抑制奖励尺度漂移和偏好语义漂移；当奖励更新过弱时，算法适度放松约束以保持奖励函数的适应能力。基于该机制，PIRO在每轮迭代中交替更新奖励函数与策略模型，在保持奖励函数动态适应性的同时提高多轮训练稳定性。

在实验层面，本文基于SLIME后训练框架，在TL;DR摘要生成和UltraFeedback通用指令跟随任务上验证了PIRO的有效性。TL;DR实验结果表明，PIRO在最佳迭代点取得$1.415 \pm 0.054$的奖励评分，高于SFT和固定奖励RLHF/PPO方法；去除RMSD近端约束后，模型性能在早期提升后出现明显退化，说明动态奖励学习需要配合奖励变化约束才能保持稳定。UltraFeedback实验结果表明，PIRO在MT-Bench八个维度上的平均得分达到$7.53$，优于SFT、DPO、RLHF(PPO)、RLHF(GRPO)以及无约束IRL基线，并在推理、STEM等维度上表现出较明显提升。两组实验共同说明，PIRO能够在不同任务和模型规模下缓解固定奖励模型带来的分布偏移和奖励过优化问题，并提升多轮偏好对齐训练的稳定性。

综上，本文从理论分析、算法设计和实验验证三个方面证明了受约束动态奖励学习在大语言模型偏好对齐中的可行性。与静态奖励模型驱动的传统RLHF流程相比，PIRO通过奖励函数与策略模型的交替更新，提高了奖励信号对当前策略分布的适应能力；通过RMSD近端约束，又在一定程度上避免了动态奖励学习带来的奖励漂移和训练振荡问题。本文工作为基于逆强化学习的大语言模型偏好对齐提供了一种新的实现路径。

\section{不足与展望}

尽管本文围绕受约束动态奖励学习进行了理论分析和实验验证，但仍存在一些不足，后续可从以下几个方面进一步改进。

首先，实验规模仍有进一步扩展空间。本文覆盖了摘要生成和通用指令跟随两个代表性任务，模型规模包含1B级和9B级设置，但相较当前大语言模型后训练实践仍然有限。后续可在更大规模模型、更长上下文任务、多轮对话、复杂推理以及工具调用等场景中进一步检验PIRO的适用性，分析RMSD近端约束在不同任务分布和策略更新强度下的稳定性。

其次，奖励建模形式仍相对单一。本文采用标量奖励函数刻画模型响应质量，并使用单一RMSD统计量约束相邻轮次奖励函数变化。然而，真实人类偏好通常同时包含事实性、安全性、帮助性、简洁性和风格一致性等多个维度，单一奖励信号难以充分表达不同目标之间的差异与冲突。后续可进一步研究多维奖励建模方法，将不同偏好维度分别纳入奖励学习和近端约束过程，从而提升奖励建模的细粒度表达能力。

再次，本文的数据来源仍以离线数据为主。实验中的高质量示范响应和偏好信号主要来自离线数据集或预先生成的数据流程，尚未充分利用真实在线交互中的用户反馈。实际应用中，用户偏好可能随任务类型、时间变化和使用群体差异而动态变化。后续可在严格隐私保护和安全审查的前提下，将PIRO与在线偏好反馈结合，使奖励函数能够在持续交互过程中更新，同时抑制在线反馈噪声、偏置样本和短期奖励诱导带来的不稳定问题。

最后，理论分析与工程效率仍有改进空间。本文的理论推导依赖有限响应空间、奖励有界和局部小步更新等假设，RMSD近端约束也是对最大逐点奖励变化的批次级近似。后续仍需进一步建立更严格的理论界，刻画采样误差、函数逼近误差、策略优化误差与奖励约束之间的关系。同时，PIRO相较固定奖励方法引入了额外的奖励更新和约束调节开销，后续可从奖励更新频率、样本复用、低秩适配和分布式训练调度等方面提升算法效率。

总体而言，受约束动态奖励学习为大语言模型偏好对齐提供了一种兼顾奖励适应性与训练稳定性的思路。未来若能在更大规模模型、更复杂反馈场景和更严格理论条件下进一步验证其有效性，该方向有望成为传统固定奖励RLHF范式的重要补充。
